Networked agent systems represent relationships in three different places, and the
three are routinely compared as one. Some structure decides \emph{who can be reached}:
which agents a requester may discover and address at all. Some decides \emph{what
crosses an edge}: whether a counterpart accepts a task and returns a result, or
exposes a store to be listed and read. The rest is \emph{asserted about an edge}:
trust scores, stated alignment, attribution, tenure---metadata an individual model
must interpret. Comparisons of open against closed agent networks vary all three at
once and attribute the difference to the regime. We separate them. Holding the model,
the runtime, the tasks and the coordination kernel fixed, we intervene on edge
formation and edge semantics as a $2\times2$ factorial, and independently seed five
relational properties as context, using a length-matched placebo to distinguish a
relational effect from the cost of the tokens carrying it. Across 144 episodes of
hidden-profile coordination scored against forced-value answer keys, a naive
public-versus-private contrast of $0.313$ decomposes into $0.232$ from edge formation
and $0.080$ from edge semantics, with the attribute axis indistinguishable from zero:
its two nominally significant contrasts, out of 56 tests per beat against a chance
floor of 2.8, reverse sign between beats of the same episodes, and the two assertions
deployed reputation systems actually implement---prior trust and attribution---are the
flattest of the set. Exposing a store buys accuracy but costs 11k tokens an episode,
so per token spent, delegating to an agent that already holds the context wins in both
arms. Before asking how internet-scale agent collectives adapt, it is worth knowing
which relational variables are interventions at all: changing the network changes the
collective, and describing the relationship does not.
